\section*{Appendix A: iFOPNG Theorem and Proof}
\label{app:a}

% ─────────────────────────────────────────────────────────────────────────────
% PREAMBLE REQUIREMENT: this appendix uses the environments
%   assumption, theorem, lemma, proposition, corollary, remark.
% If any are missing from the preamble, add (after \usepackage{amsthm}):
%   \newtheorem{lemma}{Lemma}
%   \newtheorem{proposition}{Proposition}
%   \newtheorem{corollary}{Corollary}
%   \theoremstyle{remark}\newtheorem{remark}{Remark}
% ─────────────────────────────────────────────────────────────────────────────

\subsection*{Notation}

Throughout this proof, $g \in \mathbb{R}^p$ denotes the current gradient,
$G \in \mathbb{R}^{p \times m}$ the memory matrix of $m$ historical gradient
directions (each orthonormalised at insertion via QR; the proof does not
rely on this property), and $F_{\mathrm{new}}$, $F_{\mathrm{old}}$ the
current-task and accumulated diagonal empirical Fisher tensors, with
$F_c = F_{\mathrm{new}} + F_{\mathrm{old}}$.
Diagonal Fisher tensors are treated as vectors acting by element-wise
multiplication and are identified with the corresponding diagonal matrices,
which are symmetric.
For a symmetric positive-definite $S$ we write $\|w\|_S = \sqrt{w^\top S w}$.
The projection kernel is
\begin{equation}
    A \coloneqq G^\top F_{\mathrm{old}}\, F_c^{-1}\, F_{\mathrm{old}}\, G
    \;\in\; \mathbb{R}^{m \times m}.
\end{equation}

\begin{assumption}[Regularity]
\label{ass:regularity}
$F_c$ is positive definite and the kernel $A$ is invertible.
\end{assumption}

\noindent
Both conditions hold unconditionally in the implementation, which inverts the
damped matrices $F_c + \hat\lambda I$ and $A + \hat\lambda I$ with
$\hat\lambda > 0$ (Section~\ref{sec:normalization}); the theorem below states
the exact $\lambda \to 0$ result.

\begin{theorem}[iFOPNG Update]
\label{thm:ifopng}
Under Assumption~\ref{ass:regularity}, the solution to
\begin{equation}
\label{eq:ifopng_problem}
\begin{aligned}
    \min_{u}\;\; & \bigl\| g - F_c^{-1} u \bigr\|_{F_c}^2 \\
    \mathrm{s.t.}\;\; & F_{\mathrm{old}}^{-1} u \in \operatorname{Col}(G), \\
                       & \|v\|_{F_c} = \eta,
\end{aligned}
\end{equation}
with parameter update $v = \gamma\, F_c^{-1}(g - u)$, is given by
Equation~\eqref{eq:ifopng}:
\begin{equation}
    v^* = -\eta\,
    \frac{F_c^{-1} P g}{\sqrt{(Pg)^\top F_c^{-1} P g}},
    \qquad
    P = I - F_{\mathrm{old}}\, G\, A^{-1}\, G^\top F_{\mathrm{old}}.
\end{equation}
\end{theorem}

\begin{proof}
\textbf{Step 1 (parameterise the constraint).}
The constraint $F_{\mathrm{old}}^{-1} u \in \operatorname{Col}(G)$ means
$F_{\mathrm{old}}^{-1} u = G\alpha$ for some $\alpha \in \mathbb{R}^m$,
i.e.\ $u = F_{\mathrm{old}}\, G\, \alpha$.
The optimisation over the constrained $u$ thus reduces to an unconstrained
optimisation over $\alpha \in \mathbb{R}^m$.

\textbf{Step 2 (expand the objective).}
Writing $\|w\|_{F_c}^2 = w^\top F_c\, w$ and using $F_c^{-1} F_c = I$,
\begin{equation}
    \bigl\| g - F_c^{-1} u \bigr\|_{F_c}^2
    = g^\top F_c\, g \;-\; 2\, g^\top u \;+\; u^\top F_c^{-1} u .
\end{equation}
Substituting $u = F_{\mathrm{old}}\, G\, \alpha$, the linear term becomes
$-2\,(G^\top F_{\mathrm{old}}\, g)^\top \alpha$ and the quadratic term
becomes $\alpha^\top A\, \alpha$, since
$G^\top F_{\mathrm{old}}\, F_c^{-1}\, F_{\mathrm{old}}\, G = A$.
The first term $g^\top F_c\, g$ does not depend on $\alpha$.

\textbf{Step 3 (minimise).}
The reduced objective is quadratic in $\alpha$ with Hessian $2A \succ 0$,
hence strictly convex.
Setting its gradient to zero gives the normal equations
$A\,\alpha^* = G^\top F_{\mathrm{old}}\, g$, with unique solution
\begin{equation}
\label{eq:alpha_star}
    \alpha^* = A^{-1}\, G^\top F_{\mathrm{old}}\, g .
\end{equation}

\textbf{Step 4 (back-substitute).}
Inserting $\alpha^*$ into $u = F_{\mathrm{old}}\, G\, \alpha$ gives
$u^* = F_{\mathrm{old}}\, G\, A^{-1}\, G^\top F_{\mathrm{old}}\, g$, so
\begin{equation}
    g - u^* = \bigl(I - F_{\mathrm{old}}\, G\, A^{-1}\,
              G^\top F_{\mathrm{old}}\bigr) g = P g .
\end{equation}

\textbf{Step 5 (apply the trust region).}
For $v = \gamma\, F_c^{-1} P g$,
\begin{equation}
    \|v\|_{F_c}^2
    = \gamma^2\,(Pg)^\top F_c^{-1} F_c\, F_c^{-1} (Pg)
    = \gamma^2\,(Pg)^\top F_c^{-1} (Pg).
\end{equation}
Imposing $\|v\|_{F_c} = \eta$ fixes
$|\gamma| = \eta / \sqrt{(Pg)^\top F_c^{-1} P g}$.
The negative root is selected so that with empty memory ($P = I$) the update
reduces to the natural-gradient descent step under $F_c$, matching the
\texttt{FOPNG} convention \citep{garg_fisher-orthogonal_2026}.
This yields the stated update:
\begin{equation}
    v^* = -\eta\, \frac{F_c^{-1} P g}{\sqrt{(Pg)^\top F_c^{-1} P g}}.
\end{equation}
\end{proof}

\begin{remark}[Relationship to FOPNG]
The single substitution $F_{\mathrm{new}} \to F_c$ recovers \texttt{FOPNG}:
the proof is identical with $F_c$ replaced by $F_{\mathrm{new}}$ in the
metric, the kernel $A$, and the trust region.
\texttt{iFOPNG} differs only in that the ambient metric carries the
accumulated historical curvature $F_{\mathrm{old}}$ in addition to the
current-task curvature.
\end{remark}

\begin{remark}[Inertia mechanism]
The trust-region budget splits additively across the two geometries:
\begin{equation}
    \|v\|_{F_c}^2
    = v^\top F_{\mathrm{new}}\, v + v^\top F_{\mathrm{old}}\, v
    = \|v\|_{F_{\mathrm{new}}}^2 + \|v\|_{F_{\mathrm{old}}}^2 = \eta^2 .
\end{equation}
A coordinate with large $(F_{\mathrm{old}})_i$ consumes more of the shared
budget $\eta^2$ per unit of motion and is therefore moved less.
This is the sense in which historical importance acts as inertia rather than
as a restoring force: no reference point and no penalty term enter the
objective, in contrast to EWC \citep{kirkpatrick_overcoming_2017}.
\end{remark}

\subsection*{Algorithm Statements}

Algorithm~\ref{alg:prefisher} is reproduced from
\citet{garg_fisher-orthogonal_2026} without modification.
It pre-multiplies each stored gradient by the Fisher estimate of the task it
was collected on, absorbing the metric weighting into $\tilde{G}$ at
collection time so that $\tilde{G}$ approximates $F_{\mathrm{old}}\,G$ (up to
subsequent element-wise maximum updates of $F_{\mathrm{old}}$) and the
per-iteration projection reduces to a single matrix--vector product against
$\tilde{G}$.
Empirically the two algorithms produce indistinguishable results on all
benchmarks evaluated here; \texttt{preIFOPNG} is the implementation preferred
in practice due to its reduced cost.

% ── Algorithm 1: iFOPNG ──────────────────────────────────────────────────────

\begin{algorithm}[htbp]
\caption{iFOPNG Algorithm (extends \texttt{FOPNG}~\citep{garg_fisher-orthogonal_2026}
       via combined metric $F_c$)}
\label{alg:ifopng}
\small
\setlength{\abovedisplayskip}{2pt}
\setlength{\belowdisplayskip}{2pt}
\begin{algorithmic}[1]
\Require Initial parameters $\theta_0$, task sequence $\{D_1,\ldots,D_T\}$,
              learning rate $\eta$, damping $\lambda$,
              number of stored gradients per task $k$, number of epochs $E$
\State Initialize gradient memory $G \leftarrow [\,]$.
       Train on task $D_1$ using Adam or SGD.
\State Compute $F_{\mathrm{old}} \leftarrow F(D_1)$;\quad
       Store $G \leftarrow [g_1, \ldots, g_k]$
\For{task $t = 2, 3, \ldots, T$}
       \For{epoch $e = 1, \ldots, E$}
              \State Recompute new-task Fisher:
                     ${F}_{\mathrm{new}} \leftarrow F(D_t)$;
              \State Form combined metric:
                     $F_c \leftarrow F_{\mathrm{new}} + F_{\mathrm{old}}$
              \For{all batches $(x, y)$ in $D_t$}
              \State Compute gradient $g \leftarrow \nabla_\theta \mathcal{L}_t(\theta)$
              \State Compute $v^*$ according to Theorem~\ref{thm:ifopng}
                     under metric $F_c$
              \State Update parameters: $\theta \leftarrow \theta + v^*$
              \EndFor
       \EndFor
       \State Store gradients:
              $G \leftarrow [G,\; g_{k(t-1)+1}, \ldots, g_{kt}]$
       \State Update accumulated Fisher:
              $F_{\mathrm{old}} \leftarrow
              \max\!\left(F_{\mathrm{old}},\, F_{\mathrm{new}}\right)$
\EndFor
\end{algorithmic}
\end{algorithm}

% ── Algorithm 2: preIFOPNG ───────────────────────────────

\begin{algorithm}[htbp]
\caption{preIFOPNG Algorithm~\citep{garg_fisher-orthogonal_2026},
       adapted for the iFOPNG combined metric}
\label{alg:prefisher}
\small
\setlength{\abovedisplayskip}{2pt}
\setlength{\belowdisplayskip}{2pt}
\begin{algorithmic}[1]
\Require Initial parameters $\theta_0$, task sequence $\{D_1,\ldots,D_T\}$,
              learning rate $\eta$, damping $\lambda$,
              number of stored gradients per task $k$, number of epochs $E$
\State Initialize pre-Fisher memory $\tilde{G} \leftarrow [\,]$.
       Train on task $D_1$ using Adam or SGD.
\State Compute $F_1 \leftarrow F(D_1)$;\quad
       $F_{\mathrm{old}} \leftarrow F_1$;\quad
       Store $\tilde{G} \leftarrow [F_1 \odot g_1, \ldots, F_1 \odot g_k]$
\For{task $t = 2, 3, \ldots, T$}
       \For{epoch $e = 1, \ldots, E$}
              \State Recompute new-task Fisher:
                     $F_{\mathrm{new}} \leftarrow F(D_t)$;
              \State Form combined metric:
                     $F_c \leftarrow F_{\mathrm{new}} + F_{\mathrm{old}}$
              \For{all batches $(x, y)$ in $D_t$}
              \State Compute gradient $g \leftarrow \nabla_\theta \mathcal{L}_t(\theta)$
              \State Compute $v^*$ according to Theorem~\ref{thm:ifopng}
                     using $\tilde{G}$ in place of $G$
              \State Update parameters: $\theta \leftarrow \theta + v^*$
              \EndFor
       \EndFor
       \State Compute $F_t \leftarrow F(D_t)$;\quad
              Store $\tilde{G} \leftarrow
              [\tilde{G},\; F_t \odot g_{k(t-1)+1},
              \ldots, F_t \odot g_{kt}]$
       \State Update accumulated Fisher:
              $F_{\mathrm{old}} \leftarrow
              \max\!\left(F_{\mathrm{old}},\, F_{\mathrm{new}}\right)$
\EndFor
\end{algorithmic}
\end{algorithm}